\documentclass[10pt,a4paper]{amsart}
\usepackage[margin=19mm]{geometry}
\usepackage{amsmath,amssymb,mathtools}
\usepackage[scr=rsfs,cal=euler]{mathalfa}
\usepackage{xfrac}
\usepackage[unicode,hidelinks,bookmarks=false]{hyperref}
\newcommand{\ie}{i.e.\ }
\newcommand{\cf}{cf.\ }
\newcommand{\resp}{respectively\ }
\newcommand{\Subsection}{\subsection}
\providecommand{\nobreakdash}{\nobreak\hbox{-}}
\setlength{\emergencystretch}{2em}
\multlinegap0pt
\usepackage{bm}
\usepackage{bbm}
\usepackage{dsfont}
\usepackage{xspace}
\usepackage{cancel}
\usepackage{mathtools}
\usepackage{amsthm}
\makeatletter
\@ifundefined{lemma}{\newtheorem{lemma}{Lemma}}{}
\makeatother
\title[De Rham-Higgs comparison for mixed Hodge modules in positive]{De Rham-Higgs comparison for mixed Hodge modules in positive characteristic}
\author{Zhang Zebao}
\date{Source: arXiv:2507.15175v1 (CC BY 4.0)}
\begin{document}
\maketitle

\noindent\emph{Source and licence.} De Rham-Higgs comparison for mixed Hodge modules in positive characteristic. Version arXiv:2507.15175v1 (CC BY 4.0), distributed under the Creative Commons Attribution 4.0 International licence, as recorded in the arXiv licence metadata for this exact version and shown on the abstract page https://arxiv.org/abs/2507.15175v1. Full rights notice: licenses/arxiv-2507.15175-CC-BY-4.0.txt.\par\smallskip

\noindent\textit{Editorial scope.} Counted unit: the Lemma whose source label is \texttt{None} in the article (source0.tex lines 1407--1409, source label \texttt{None}), delivered together with the article's own complete original proof of it (source0.tex lines 1410--1528, one \texttt{proof} environment of 119 lines). No mathematics is written, repaired or re-derived: statement and proof are the article's own text line for line. Required context retained uncounted: none. Every definition, display and label of those lines is retained unchanged. Omitted from this package and not redistributed: 1--1406, 1529--6244. Nothing inside a retained excerpt is omitted or altered: every retained body is delivered exactly as the article's lines are written, so no omission receipt of any kind accompanies this package. Citations: the retained excerpts carry no citation command at all, so no bibliography entry of the article is transcribed and no reference is redistributed. Reference adaptation: none was needed and none was made; every reference inside the delivered unit resolves inside it, so no unresolved reference or citation is produced. Printed slips kept literal: no printed word, symbol, hypothesis or number of the counted statement or of its proof was altered. Layout and class: the article's own class is \texttt{article} with packages including geometry, amsmath, amsthm, amssymb, amscd, bm, mathrsfs, dsfont, cases, bbding, hyperref, xy, upgreek, xcolor; the wrapper is the lane's pinned \texttt{amsart} editorial wrapper at 10pt on A4, which restates the theorem-like environments the retained text needs and adds the article's own macro definitions that the retained text needs (none), with extra wrapper packages (bm, bbm, dsfont, xspace, cancel, mathtools). The delivered numbering is the unit's own. Assets: no retained span contains a figure environment, a graphics-inclusion command, a PDF-page inclusion, a table or a TikZ picture; no image file is copied into this package, the package declares no asset dependency and no image file is redistributed with it. Licence: version arXiv:2507.15175v1 of this article is distributed under the Creative Commons Attribution 4.0 International licence, as recorded in the arXiv licence metadata for this exact version and shown on the abstract page https://arxiv.org/abs/2507.15175v1. A full rights notice is delivered at licenses/arxiv-2507.15175-CC-BY-4.0.txt. Characters: every retained excerpt is pure ASCII, so the delivered engine, XeLaTeX with its default Latin Modern text font, reports no missing character.
\begin{lemma}
The complexes $\star(\hat E,\hat\theta-df\wedge)$ and $\star(\hat E,\Theta)$ are isomorphic.   
\end{lemma}

\begin{proof}
The proof contains three steps: the first step is to prove this lemma for $\star=\Omega^*$; the second step is for $\Omega^*_{int},W_i\Omega^*$; the last step is for $\star=\Omega^*_{g^{-1}}$.

\emph{Step 1}. Recall that the assumption (S1) gives rise to a system of coordinates $t_1,\cdots,t_n$ on $X/k$ such that $D=(t_1\cdots t_r=0)$ and $g=t_1\cdots t_s$. Write
\begin{eqnarray}\label{Theta cor}
\hat\theta=\sum_{i=1}^r\hat\theta_i\otimes d\log t_i+\sum_{i=r+1}^n\hat\theta_i\otimes dt_i,~\Theta=\sum_{i=1}^r\Theta_i\otimes d\log t_i+\sum_{i=r+1}^n\Theta_i\otimes dt_i
\end{eqnarray}
and the K\"ahler differential 
$$d=\sum_{i=1}^rt_i\partial_i\otimes d\log t_i+\sum_{i=r+1}^n\partial_i\otimes dt_i.$$ 
Clearly, we have natural isomorphisms
$$
\Omega^*(\hat E,\hat\theta-df\wedge)\cong\mathrm{Kos}(\hat E;\hat\theta_1-t_1\partial_1f,\cdots,\hat\theta_r-t_r\partial_rf,\hat\theta_{r+1}-\partial_{r+1}f,\cdots,\hat\theta_n-\partial_nf)
$$
and
$$
\Omega^*(\hat E,\Theta)\cong\mathrm{Kos}(\hat E;\Theta_1,\cdots,\Theta_n).
$$
To show $\Omega^*(\hat E,\hat\theta-df\wedge)$ is isomorphic to $\Omega^*(\hat E,\Theta)$, it suffices to check that 
$$
\Theta_i=\left\{
\begin{matrix}
  \vartheta_i(\hat\theta_i-t_i\partial_if),&1\leq i\leq r;\\
\vartheta_i(\hat\theta_i-\partial_if),&r+1\leq r\leq n,
\end{matrix}
\right.
$$
where each $\vartheta_i$ is an $\mathcal O_{\mathfrak X}$-linear automorphism of $\hat E$. In fact, we may construct an isomorphism of complexes 
$$
\psi:\Omega^*(\hat E,\hat\theta-df\wedge)\to\Omega^*(\hat E,\Theta) 
$$
as follows: set 
\begin{eqnarray}\label{omega}
\omega_1:=d\log t_1,\cdots,\omega_r:=d\log t_r,~\omega_{r+1}:=dt_{r+1},\cdots,\omega_n:=dt_n
\end{eqnarray}
and put
$$
\psi(\hat e):=\hat e,~\psi(\hat e\otimes\omega_{i_1}\wedge\cdots\wedge\omega_{i_q}):=\vartheta_{i_1}\cdots\vartheta_{i_q}(\hat e)\otimes\omega_{i_1}\wedge\cdots\wedge\omega_{i_q},~\hat e\in\hat E.
$$

Next, let us construct $\vartheta_i$. Using the facts
$$
t_q\partial_qf_{i_1,\cdots,i_n}=i_qf_{i_1,\cdots,i_n};~i_q=i_q^p\mod p;~x^p+y^p=(x+y)^p,~x,y\in\mathcal O_{\mathfrak X},
$$
we get
\begin{eqnarray*}
\sum_{0\leq i_1,\cdots,i_n<p}\sum_{j=0}^{\infty}f_{i_1\cdots,i_n}^{p^j-1}df_{i_1,\cdots,i_n}=\sum_{0\leq i_1,\cdots,i_n<p}\sum_{j=0}^{\infty}\sum_{q=1}^nf_{i_1\cdots,i_n}^{p^j-1}f_{i_1,\cdots,i_n}t_q\partial_qf_{i_1,\cdots,i_n}d\log t_q\\
=\sum_{0\leq i_1,\cdots,i_n<p}\sum_{j=0}^{\infty}\sum_{q=1}^n(i_qf_{i_1\cdots,i_n})^{p^j}d\log t_q=\sum_{q=1}^n\sum_{j=0}^{\infty}(t_q\partial_qf)^{p^j}d\log t_q.~~~~~~~~~~
\end{eqnarray*}
It follows that
$$
\Theta_i=\hat\theta_i-\sum_{q=1}^n\sum_{j=0}^{\infty}(t_q\partial_qf)^{p^j},~1\leq i\leq r;~t_i\Theta_i=t_i\hat\theta_i-\sum_{q=1}^n\sum_{j=0}^{\infty}.(t_q\partial_qf)^{p^j},~r+1\leq i\leq n.
$$
Consequently, for $1\leq i\leq r$, we may set
\begin{eqnarray*}
\vartheta_i:=\frac{\hat\theta_i-\sum_{j=0}^\infty(t_i\partial_i f)^{p^j}}{\hat\theta_i-t_i\partial_i f}=1-\frac{\sum_{j=1}^\infty(t_i\partial_i f)^{p^j}}{\hat\theta_i-t_i\partial_i f}~~~~~~\\
=1+\sum_{j=1}^\infty\sum_{q=0}^{p-1}(t_i\partial_i f)^{p^j-1}(\frac{\hat\theta_i}{t_i\partial_if})^q=1+\sum_{j=1}^\infty\sum_{q=0}^{p-1}(t_i\partial_i f)^{p^j-q-1}\hat\theta_i^q.
\end{eqnarray*}
Here the third equality follows from the Taylor expansion
$\frac{1}{1-x}=\sum_{i=0}^\infty x^i$
and the nilpotent condition $\hat\theta_i^p=0$. Since each term $(t_i\partial_i f)^{p^j-s-1}\hat\theta_i^s$ is a topological nilpotent endomorphism of $\hat E$, we conclude that $\vartheta_i$ is an automorphism of $\hat E$. Similarly, for $r+1\leq i\leq n$, we may set
$$
\vartheta_i:=1+\sum_{j=1}^\infty\sum_{q=0}^{p-1}(t_i\partial_i f)^{p^j-q-1}(t_i\hat\theta_i)^q.
$$
Clearly, it is an automorphism of $\hat E$ again.This completes the construction of $\vartheta_i$.

\emph{Step 2}. Using the following easily checked facts
$$
\psi(\Omega_{int}^j(\hat E,\hat\theta-df\wedge))=\Omega_{int}^j(\hat E,\Theta),~\psi(W_i\Omega^j(\hat E,\hat\theta-df\wedge))=W_i\Omega^j(\hat E,\Theta),
$$
we deduce that the isomorphism $\psi$ restricts to isomorphisms of subcomplexes
$$
\psi_{int}:\Omega_{int}^*(\hat E,\hat\theta-df\wedge)\to\Omega_{int}^*(\hat E,\Theta),~W_i\psi:W_i\Omega^*(\hat E,\hat\theta-df\wedge)\to W_i\Omega^*(\hat E,\Theta).
$$ 

\emph{Step 3}. Unfortunately, the isomorphism $\psi$ constructed in \emph{step 1} may not restricts to an isomorphism of complexes 
$$\psi_{g^{-1}}:\Omega^*_{g^{-1}}(\hat E,\hat\theta-df\wedge)\to\Omega^*_{g^{-1}}(\hat E,\Theta).$$
Nevertheless, we may construct the desired isomorphism $\psi_{g^{-1}}$ by using a similar idea as $\psi_{int}$ or $W_i\psi$. Write
$$
\hat\theta=\hat\theta'_1\otimes d\log g+\sum_{i=2}^r\hat\theta'_2\otimes d\log t_i+\sum_{i=r+1}^n\hat \theta'_i\otimes dt_i.
$$
and
$$
d=t_1\partial_1\otimes d\log g+\sum_{i=2}^s(t_i\partial_i-t_1\partial_1)\otimes d\log t_i+\sum_{i=s+1}^rt_i\partial_i\otimes d\log t_i+\sum_{i=r+1}^n\partial_i\otimes dt_i.
$$
it is easy to see that
$$
\hat\theta_1'=\hat\theta_1,~\hat\theta_2'=\hat\theta_2-\hat\theta_1,\cdots,\hat\theta_s'=\hat\theta_s-\hat\theta_1,~\hat\theta_{s+1}'=\hat\theta_{s+1},\cdots,\hat\theta_n'=\hat\theta_n
$$
and 
$$
df=t_1\partial_1f\otimes d\log g+\sum_{i=2}^s(t_i\partial_i-t_1\partial_1)f\otimes d\log t_i+\sum_{i=s+1}^rt_i\partial_if\otimes d\log t_i+\sum_{i=r+1}^n\partial_if\otimes dt_i.
$$
For our purpose, we construct another isomorphism of complexes 
$$
\psi':\Omega^*(\hat E,\hat\theta-df\wedge)\to\Omega^*(\hat E,\Theta) 
$$
as follows: set
$$
\omega'_1:=d\log g,~\omega'_2:=d\log t_2,\cdots,\omega'_r:=d\log t_r,~\omega'_{r+1}:=dt_{r+1},\cdots,\omega'_n:=dt_n
$$
and put
$$
\psi'(\hat e)=\hat e,~\psi'(\hat e\otimes\omega_{i_1}\wedge\cdots\omega_{i_q})=\vartheta'_{i_1}\cdots\vartheta'_{i_q}(\hat e)\otimes\omega'_{i_1}\wedge\cdots\wedge\omega'_{i_q}.
$$
Here 
$$
\vartheta_i':=\left\{
\begin{matrix}
1+\sum_{j=1}^\infty\sum_{q=0}^{p-1}((t_i\partial_i-t_1\partial_1) f)^{p^j-q-1}\hat\theta_i^q,&2\leq i\leq s,\\
\vartheta_i,&\mathrm{otherwise}.
\end{matrix}
\right.
$$
By the very construction of $\Omega^*_{g^{-1}}$, it is easy to see that $\psi'$ restricts to an isomorphism of subcomplexes 
$$
\psi'_{g^{-1}}:\Omega^*_{g^{-1}}(\hat E,\hat\theta-df\wedge)\to\Omega^*_{g^{-1}}(\hat E,\Theta).
$$
This completes the proof. 
\end{proof}

\end{document}
